\documentclass[10pt,conference]{IEEEtran}
% ICSE'27 instructions: \documentclass[10pt,conference]{IEEEtran}, no compsoc options.
\usepackage{cite}
\usepackage{amsmath,amssymb}
\usepackage{graphicx}
\usepackage{booktabs}
\usepackage{multirow}
\usepackage{xcolor}
\usepackage{url}
\usepackage[hidelinks]{hyperref}

% AUTO-GENERATED by scripts/paper_numbers.py -- do not edit.
\newcommand{\BaseGflopsM}{2.67}
\newcommand{\BaseMeanFone}{0.910}
\newcommand{\BaseScoreLocal}{0.730}
\newcommand{\BaseScoreOfficial}{0.729}
\newcommand{\BaseSortedGflopsM}{1.17}
\newcommand{\BaseSortedScore}{0.852}
\newcommand{\CloneAgreePct}{98}
\newcommand{\CloneTestPct}{18}
\newcommand{\CostModelErr}{0.04}
\newcommand{\CostModelN}{13}
\newcommand{\MatNoTagPct}{12}
\newcommand{\MatPrev}{1.65}
\newcommand{\PadEff}{43}
\newcommand{\RuleMatFone}{0.935}
\newcommand{\RuleMatP}{1.000}
\newcommand{\RuleMatR}{0.877}
\newcommand{\StudentLayers}{3}
\newcommand{\TestN}{30{,}655}
\newcommand{\TokOurs}{154}
\newcommand{\TokRaw}{235}
\newcommand{\TrainN}{214{,}586}
\newcommand{\TruncOurs}{0.5}
\newcommand{\TruncRaw}{10.6}
\newcommand{\ValN}{21{,}243}
\newcommand{\VulnPrev}{32.8}
\newcommand{\VulnQFour}{60.3}
\newcommand{\VulnQOne}{8.5}
 % generated by scripts/paper_numbers.py from experiment outputs; never edit by hand

\begin{document}

\title{AEGIS: Efficiency-Guided Multitask Learning for Joint Vulnerability and Technical-Debt Detection}

\author{\IEEEauthorblockN{Awaiz Ahmed}
\IEEEauthorblockA{\textit{Affiliation (to be completed)}\\
email (to be completed)}}

\maketitle

\begin{abstract}
% Written last, from measured results only.
\end{abstract}

\begin{IEEEkeywords}
vulnerability detection, self-admitted technical debt, multitask learning, knowledge distillation, efficient inference
\end{IEEEkeywords}

\section{Introduction}
The NLBSE'27 tool competition asks for a single model that reads a C/C++ function and predicts two
labels: whether the function is associated with a weakness or vulnerability, and whether it carries a
self-admitted technical debt (SATD) annotation~\cite{nlbse2027}. Submissions are ranked by a score
that weights mean F1 by 0.6, and inference runtime and GFLOPs, measured on a Colab T4 over the whole
test split, by 0.2 each. The organizers' CodeBERT~\cite{feng2020codebert} baseline reaches a mean F1 of
0.910 but uses almost the entire GFLOPs budget.

AEGIS treats the ranking score itself as the design objective. We (i)~derive the exchange rates the
score implies between F1, runtime and FLOPs and an analytic cost model that matches the official
profiler count, (ii)~shrink the processed input before touching the model (whitespace normalization, a
compact tokenizer, length-sorted dynamic padding), and (iii)~spend the remaining budget on a shallow
student distilled from a full-depth multitask teacher, with per-task decision thresholds calibrated on
held-out training data.

\section{Task, Data and Scoring}
\label{sec:task}
% Dataset: MADE-WIC derived (OSPR, Big-Vul, Devign) \cite{mock2024madewic}; label semantics (MAT);
% splits; audit findings: MAT keyword ceiling, length confound, clone overlap.

% Scoring: formula, measurement loop, exchange rates (0.01 F1 = 0.006 = 90k GFLOPs = 30 s),
% profiler coverage (linear layers only).

\section{The AEGIS Approach}
\label{sec:method}
% Input compaction; dynamic padding; teacher; layer-truncated student; multitask distillation loss;
% threshold calibration; final checkpoint is a standard RobertaForSequenceClassification.

\section{Experimental Setup}
\label{sec:setup}
% Internal validation split; hardware; hyperparameters; seeds; measurement protocol.

\section{Results}
\label{sec:results}
% Main table (public test, official loop); frontier figure; ablations; error analysis.

\section{Limitations}
\label{sec:limitations}

\section{Conclusion}
\label{sec:conclusion}

\section*{Acknowledgment}
% IEEE policy (quoted in the ICSE'27 instructions): use of AI-generated content must be disclosed here.

\bibliographystyle{IEEEtran}
\bibliography{references}

\end{document}
